\documentclass[10pt]{article}
\usepackage[margin=0.8in]{geometry}
\usepackage{amsmath,amssymb,amsthm,booktabs,microtype,xcolor,hyperref}
\hypersetup{colorlinks=true,linkcolor=blue!45!black,urlcolor=blue!45!black}
\newtheorem{prop}{Proposition}
\newcommand{\Z}{\mathbb Z}\newcommand{\Q}{\mathbb Q}
\newcommand{\TL}{\mathrm{TL}}\newcommand{\id}{\mathbf 1}
\setlength{\parindent}{0pt}\setlength{\parskip}{5pt}
\begin{document}
{\LARGE A bounded attack on Jones unknot detection}\\[3pt]
{\large Exact results and the remaining generic obstruction}\\
31 October 2026 \hfill Research note; no claim of novelty

\textbf{Outcome.} This pass produced neither a generic proof nor a nontrivial knot
with Jones polynomial \(1\). It tested general mechanisms rather than extending
the earlier EKT-family computations. The results below delimit possible
approaches; they do not resolve the conjecture.

\textbf{1. No fixed Taylor cutoff works, even on three strands.}
Normalize \(V_U=1\), and set \([x,y]=xyx^{-1}y^{-1}\).
In \(B_3\), define
\[
 a=\sigma_1^2,\quad b=\sigma_2^2,\quad
 c_1=a,\quad c_{r+1}=[c_r,b],\quad
 K_r=\widehat{c_r\sigma_1\sigma_2}.
\]
Each \(c_r\) is pure, so the closure is a knot: its permutation is a three-cycle.
\begin{prop}
For every \(m\geq1\),
\[
 V_{K_{2m}}(t)=1-3\,4^{m-1}(t-1)^{2m}
                  +O((t-1)^{2m+1}).
\]
Thus \(K_{2m}\) is nontrivial, but its Jones derivatives through order \(2m-1\)
agree with those of the unknot.
\end{prop}
\begin{proof}
Put \(t=A^{-4}\), \(h=A-1\), and \(\delta=-A^2-A^{-2}\).
The writhe-normalized Temperley--Lieb representation has
\[
 R_i=-A^{-2}\id-A^{-4}e_i,\qquad
 e_i^2=\delta e_i,\quad e_1e_2e_1=e_1,\quad e_2e_1e_2=e_2.
\]
Direct multiplication gives
\[
 R_i^2=A^{-4}\id+(A^{-6}-A^{-10})e_i
       =\id+h(-4\id+4e_i)+O(h^2).
\]
The formal-series identity
\([\id+O(h^r),\id+O(h)]=\id+O(h^{r+1})\)
implies \(R(c_r)=\id+O(h^r)\). For \(r\geq2\), write
\(R(c_r)=\id+h^r C_r+O(h^{r+1})\). In the specialized algebra \(\delta=-2\),
\[
 C_2=16[e_1,e_2],\qquad C_{r+1}=4[C_r,e_2].
\]
For \(D(x)=[x,e_2]\), the relation \(e_2^2=\delta e_2\) gives
\(D^3=\delta^2D=4D\). Since every \(C_r\), \(r\geq2\), lies in the image of \(D\),
we obtain \(C_{r+2}=64C_r\).

Let \(\ell\) be braid closure with one circle valued \(1\).
At \(\delta=-2\),
\[
 \ell(\id)=4,\quad\ell(e_1)=\ell(e_2)=-2,\quad
 \ell(e_1e_2)=\ell(e_2e_1)=1.
\]
The baseline operator at \(A=1\) is
\(R(\sigma_1\sigma_2)=(\id+e_1)(\id+e_2)\).
The TL relations give
\[
 \ell([e_1,e_2](\id+e_1)(\id+e_2))=1-\delta^2=-3.
\]
Since the baseline closure is the unknot, the first nonzero coefficient of
\(V_{K_{2m}}(A^{-4})-1\) is \(-48\,64^{m-1}h^{2m}\).
The local change of variable \(t-1=-4h+O(h^2)\) gives the displayed result.
\end{proof}
For any fixed \(N\), choosing \(2m>N\) defeats a test based only on derivatives
through order \(N\). This is an explicit instance of the familiar finite-type
limitation. It does \emph{not} exclude complexity-dependent bounds or arguments
using the full Laurent polynomial.

Exact computations checked \(r=1,\ldots,8\). For \(r=2,4,6,8\), the first nonzero
\(t-1\) coefficients were respectively \(-3,-12,-48,-192\), in exactly degrees
\(2,4,6,8\). The displayed braid-word lengths were \(10,34,130,514\); these are
not claims about minimal crossing numbers. The unbounded assertion rests on
the proof, not extrapolation.

\newpage
\textbf{2. Complete classification of two-strand skein units.}

Let \(R=\Z[A,A^{-1}]\), \(\delta=-A^2-A^{-2}\), and
\[
 S=R[e]/(e^2-\delta e).
\]
Under horizontal tangle concatenation, the zero tangle represents \(\id\), the
infinity tangle represents \(e\), and a positive integer crossing represents
\(b=A^{-1}+Ae\). This horizontal convention differs from the vertical convention
on page 1.

\begin{prop}
Every unit of \(S\) has the form
\[
 \eta A^s b^k,\qquad \eta\in\{1,-1\},\quad s,k\in\Z.
\]
A product of tangle brackets can equal \(\id\) only if each factor already has
the bracket of an integer tangle up to a signed Laurent monomial.
\end{prop}
\begin{proof}
The characters of \(f+ge\) are \(f\) and \(f+\delta g\).
Both must be units if \(f+ge\) is a unit. Conversely, if both are units,
\[
 (f+ge)^{-1}=\frac{(f+\delta g)-ge}{f(f+\delta g)}.
\]
Write \(f=\eta A^m\), \(f+\delta g=\epsilon A^n\).
Their difference is divisible by \(\delta\) precisely when the monomials agree
modulo \(A^4+1\). Since \(A^4=-1\), with basis \(1,A,A^2,A^3\), this says
\[
 n-m=4k,\qquad \epsilon=\eta(-1)^k.
\]
The characters of \(b\) are \(A^{-1}\) and \(-A^3\), so \(f+ge\) and
\(\eta A^{m+k}b^k\) have the same characters. They are equal since \(\delta\)
is not a zero divisor. Negative powers are allowed:
\(b^{-1}=A+A^{-1}e\).
\end{proof}

For an \emph{actual} tangle with unit bracket, composing with an integer tangle
leaves a scalar bracket. Its denominator closure has a monomial bracket.
The identities \(V_L(1)=(-2)^{\#L-1}\) and \(V_K'(1)=0\) force that closure to
be a knot with normalized Jones polynomial \(1\). The scalar is therefore
a framing unit \((-A^3)^w\).

The missing step is geometric: must that actual tangle be the framed integer
tangle whose bracket it shares? Sikora's Theorem 9.2 makes this implication
equivalent to the Jones unknot conjecture [1]. The algebraic classification
does not supply the topological conclusion.

\textbf{Why composition with caps breaks a proposed induction.}
An explicit nonrational algebraic tangle is
\[
 T=\langle-1\rangle+\langle1/2\rangle+\langle1/3\rangle.
\]
Its bracket \(f+ge\) has
\[
 \begin{split}
 f&=A^2-2A^{-2}+2A^{-6}-A^{-10},\\
 g&=A^4+A^{-8}-A^{-12},\qquad \delta f+g=1.
 \end{split}
\]
At \(A=e^{\pi i/4}\), \(f=6i,g=1\), giving fraction \(-1/6\).
That rational tangle has monomial \(g\) up to framing, unlike \(T\).
Its six-crossing numerator independently simplifies to the unknot.
If \(Q\) is the rotation of \(T\), then
\[
 [Q]=g+fe,\qquad e[Q]=(g+\delta f)e=e.
\]
Thus a cap absorbs a nonunit tangle. An induction that allows cups, caps, or
changes in boundary strand count cannot simply apply the unit-factor result.

\textbf{Higher strand counts.}
The factor restriction extends to every \(\TL_n\): determinants of left
multiplication in its finite free \(R\)-module show that every factor of a unit
is a unit. The simple classification does not extend. In \(\TL_3\),
\[
 N=\delta e_1e_2-e_1,\quad N^2=0,\quad
 (\id+N)^{-1}=\id-N.
\]
This formal unit is not a signed-monomial multiple of a braid image. At \(A=i\),
where \(\delta=2\), braid images are unitary for \(e_i^*=e_i\), whereas
\[
 (\id+N)^*(\id+N)-\id
 =8e_2-2(e_1e_2+e_2e_1)\ne0.
\]
No realization of this formal unit by a single actual tangle is asserted.

\newpage
\textbf{3. Other generic mechanisms and their exact stopping points.}

\textbf{Branched-cover homology.}
For every determinant-one knot, the double branched cover is an integral homology
sphere. At any crossing, the two resolutions have first homology
\(\Z/n\) and \(\Z/(n+1)\), with a suitable ordering and integer \(n\);
the switched crossing gives \(\Z/(2n+1)\), where \(\Z/0=\Z\).
These follow from distance-one filling slopes on a knot exterior in that homology
sphere. Cyclic homology and consecutive smoothing determinants are thus automatic,
rather than new evidence for unknottedness.

All 44 resolutions of the standard Kinoshita--Terasaka and Conway diagrams passed
these tests by integral Fox matrices and Smith forms. Each knot has an unknotting
crossing whose oriented smoothing has \(H_1=\Z\), but Jones polynomial
\[
 x^9-x^7+x^5-x^3-x-x^{-1}-x^{-3}+x^{-5}-x^{-7}+x^{-9},
 \qquad x=t^{1/2},
\]
instead of the unlink value \(-x-x^{-1}\).
Even a band to an actual unknot plus this homology condition does not force
the smoothing to be the unlink.

\textbf{Braid representations and positivity.}
In \(B_2\), \(\sigma\) and \(\sigma^{-1}\) both close to the unknot, although
their normalized operators are different. Appending the same context
\(\sigma^2\) produces a trefoil and an unknot. Hence equality of Jones
polynomials is not a local operator replacement rule.

There is a precise all-strand positivity lemma. Normalize the Markov trace by
\(\tau(\id)=1\); then on \(n\) strands
\[
 V_{\widehat\beta}=\delta^{n-1}\tau(R(\beta)).
\]
For \(A\) on the unit circle sufficiently close to \(i\), the fixed
finite-dimensional TL trace is positive and faithful, and braid images are
unitary. If \(V_{\widehat\beta}=\delta^{n-1}\), then
\[
 \tau((R(\beta)-\id)^*(R(\beta)-\id))=0,
\]
so \(R(\beta)=\id\). Holding on an interval, this is a Laurent identity.
The character \(e_i\mapsto0\) forces writhe zero.
Thus the \(n\)-component unlink value of an \(n\)-braid forces its entire
generic TL image to be the identity.

For a \emph{knot} with \(V=1\), however, the trace is \(\delta^{1-n}\), strictly
inside the unit disk for \(n>1\). The equality case needed for this argument is
absent. A 2026 preprint claims faithfulness of the four-strand Burau and Jones
representations [4]; its source and statement were checked, not the entire proof.
Even a valid faithfulness theorem does not close the trace gap. Ito's theorem [3]
makes a nontrivial representation kernel a sufficient route to knot counterexamples,
but does not give the converse.

\textbf{Categorification.}
Reduced rational Khovanov homology detects the unknot by rank one [2].
If each quantum grading occupied only one homological parity, total rank would
equal the sum of absolute Jones coefficients, and \(V=1\) would finish the proof.
A nontrivial Jones-trivial knot therefore requires cancellation between the
two parities in some quantum grading.

The formal unreduced series
\[
 P(u,q)=q+q^{-1}+q^{-1}(1+u)(1+uq^4)
\]
has nonnegative ranks, total rank six, Euler characteristic
\(P(-1,q)=q+q^{-1}\), and a degree-\((1,4)\) differential leaving precisely
\((0,\pm1)\): pair the added \((0,-1)\) with \((1,3)\), and \((1,-1)\) with
\((2,3)\). Thus positivity of dimensions, Euler characteristic, and a Lee-shaped
endpoint alone do not force unknot homology. This abstract model is \emph{not}
claimed to arise from a knot or satisfy every Khovanov constraint.

\textbf{Status.}
No generic proof or counterexample was obtained. What remains is a restriction
on topologically realizable cancellation strong enough to turn Jones value \(1\)
into geometric triviality. None of the scalar, unit, homology, or formal homology
arguments above supplies that missing implication.

\newpage
\textbf{Verification and primary references.}

The full proofs of Propositions 1 and 2 were independently audited. The first
was checked against exact Laurent computations through \(r=8\). The second was
checked with 2,754 exact unit/inverse computations, including negative twists
and scalar factors. The tangle absorption example was checked algebraically and
by a separate diagram calculation. The branched-cover tests used all 44 crossing
resolutions, and the stated smoothing Jones polynomial was calculated using two
independent Regina algorithms.

These finite checks validate conventions and implementations; they are not the
basis for unbounded conclusions. The accompanying evidence archive contains
the working proofs, exact data, verifier scripts, and this document's source.
The basic algebra checks use Python's standard library. The optional
diagram-based replay requires Regina and Spherogram.

For the finite-order construction, a separate exact Burau recurrence additionally
proves every \(K_r\), including odd \(r\), is nontrivial. At \(t=-1\), use
\[
 A=\begin{pmatrix}1&2\\0&1\end{pmatrix},\quad
 B=\begin{pmatrix}1&0\\-2&1\end{pmatrix},\quad
 R_0=\begin{pmatrix}0&1\\-1&1\end{pmatrix},\qquad
 C_{r+1}=C_rBC_r^{-1}B^{-1}.
\]
If \(b_r,d_r\) are the top-right and bottom-right entries of \(C_r\), then
\[
 b_{r+1}=2b_r^2,\qquad d_{r+1}=1+2b_rd_r,\qquad
 b_r/2\leq d_r<b_r.
\]
The last inequalities follow by induction from \(b_1=2,d_1=1\).
Consequently
\[
 \operatorname{tr}(C_{r+1}R_0)
 =3-2d_r^2+6b_rd_r-2b_r^2
 \geq3+b_r^2/2\geq5.
\]
The Burau--Alexander formula gives
\(\det(K_r)=|2-\operatorname{tr}(C_rR_0)|>1\);
the base case is determinant \(3\). This fact is not needed for the even-order
proof on page 1.

[1] Adam S. Sikora, \emph{Tangle equations, the Jones conjecture, slopes of surfaces
in tangle complements, and \(q\)-deformed rationals}, Canadian Journal of Mathematics
76 (2024), 707--727, Theorem 9.2.
\url{https://doi.org/10.4153/S0008414X23000755}

[2] P. B. Kronheimer and T. S. Mrowka, \emph{Khovanov homology is an unknot-detector},
Publications Math\'ematiques de l'IH\'ES 113 (2011), 97--208.
\url{https://arxiv.org/abs/1005.4346}

[3] Tetsuya Ito, \emph{A kernel of a braid group representation yields a knot with
trivial knot polynomials}, Mathematische Zeitschrift 280 (2015), 347--353.
\url{https://arxiv.org/abs/1402.2028}

[4] Vasudha Bharathram, Joan S. Birman, and Tara E. Brendle,
\emph{The Burau representation is faithful for \(n=4\)}, v2, 14 September 2026.
This pass did not audit the complete proof.
\url{https://arxiv.org/abs/2607.05283}

[5] Stephen Bigelow, \emph{Does the Jones polynomial detect the unknot?}
\url{https://arxiv.org/abs/math/0012086}

[6] Dror Bar-Natan, \emph{On Khovanov's categorification of the Jones polynomial},
Algebraic \& Geometric Topology 2 (2002), 337--370.
\url{https://msp.org/agt/2002/2-1/agt-v2-n1-p16-p.pdf}

[7] Sophie Morier-Genoud, Valentin Ovsienko, and Alexander P. Veselov,
\emph{Burau representation of braid groups and \(q\)-rationals},
International Mathematics Research Notices 2024, no. 10, 8618--8627.
\url{https://doi.org/10.1093/imrn/rnad318}

\textbf{Interpretation.} These are restrictions on possible approaches, not a
resolution or a novelty claim. In particular, the explicit three-strand family
rules out a fixed Taylor cutoff; its full Jones polynomial is nontrivial.
\end{document}
